\documentclass[12pt,a4paper]{article}

\usepackage[T1]{fontenc}
\usepackage[utf8]{inputenc}
\usepackage[brazil]{babel}
\usepackage{lmodern}
\usepackage{geometry}
\usepackage{hyperref}
\usepackage{enumitem}
\usepackage{array}

\geometry{margin=2.5cm}
\hypersetup{
  colorlinks=true,
  linkcolor=blue,
  urlcolor=blue
}

\title{Plano de Aula: Aula 4 -- Engenharia de Requisitos}
\date{}

\begin{document}
\maketitle

\noindent
\textbf{Disciplina:} Engenharia de Software (Curso Técnico em Desenvolvimento de Sistemas)\\
\textbf{Duração:} 2 horas (120 minutos)\\
\textbf{Modalidade:} Presencial\\
\textbf{Materiais da aula:} \href{aula-04.html}{slides (HTML)} \quad \href{aula-04.pdf}{ementa (PDF)}\\
\textbf{Livro-base (referência):} \href{../99-Materiais/Engenharia%20de%20software%20projetos%20e%20processos%20by%20Wilson%20de%20P%C3%A1dua%20Paula%20Filho%20.epub.pdf}{Engenharia de software: projetos e processos (PDF)}

\section{Competências e Habilidades}
\textbf{Competências}
\begin{itemize}[leftmargin=*,itemsep=0.25em]
  \item Compreender o processo de engenharia de requisitos e seu impacto em prazo, custo e qualidade.
  \item Aplicar princípios de engenharia de software para reduzir retrabalho por requisitos mal definidos.
\end{itemize}

\textbf{Habilidades}
\begin{itemize}[leftmargin=*,itemsep=0.25em]
  \item Elicitar requisitos a partir de um problema e stakeholders.
  \item Diferenciar requisitos funcionais e não funcionais.
  \item Escrever requisitos testáveis com critérios de aceitação.
\end{itemize}

\section{Objetivos da Aula}
\begin{itemize}[leftmargin=*,itemsep=0.25em]
  \item Entender o fluxo de requisitos: elicitação, análise, especificação e validação.
  \item Classificar requisitos (funcionais e não funcionais) com exemplos.
  \item Produzir histórias de usuário e critérios de aceitação.
  \item Identificar problemas típicos: ambiguidade, conflitos, escopo implícito.
\end{itemize}

\section{Assuntos Abordados no Dia}
\begin{itemize}[leftmargin=*,itemsep=0.25em]
  \item Conceitos e importância de requisitos.
  \item Tipos de requisitos e características de bons requisitos.
  \item Técnicas e saídas: backlog, histórias de usuário, casos de uso.
  \item Validação e verificação: critérios de aceitação e ligação com testes.
\end{itemize}

\section{Metodologia}
Exposição dialogada com exemplos curtos e discussão orientada. Atividade prática em grupo para escrever requisitos funcionais e não funcionais e definir critérios de aceitação. O professor revisa com a turma e demonstra como transformar aceitação em casos de teste.

\section{Materiais e Recursos}
\begin{itemize}[leftmargin=*,itemsep=0.25em]
  \item Quadro e marcadores.
  \item Slides: \href{aula-04.html}{aula-04.html} e ementa: \href{aula-04.pdf}{aula-04.pdf}.
  \item Livro-base (PDF): \href{../99-Materiais/Engenharia%20de%20software%20projetos%20e%20processos%20by%20Wilson%20de%20P%C3%A1dua%20Paula%20Filho%20.epub.pdf}{Engenharia de software: projetos e processos}.
  \item Folha A4 para grupo e canetas.
  \item Vídeos complementares (links nos slides).
\end{itemize}

\section{Cronograma e Descrição Detalhada (120 Minutos)}
\subsection*{Parte 1: Abertura e motivação (10 min)}
\begin{itemize}[leftmargin=*,itemsep=0.35em]
  \item Relembrar o caso do estoque (Aula 1) e conectar com requisitos.
  \item Apresentar objetivos e o formato da atividade.
\end{itemize}

\subsection*{Parte 2: Conceitos e tipos de requisitos (30 min)}
\begin{itemize}[leftmargin=*,itemsep=0.35em]
  \item Definir requisitos e stakeholders.
  \item Requisitos funcionais x não funcionais: exemplos do ``app organizador''.
  \item Características de bons requisitos: clareza, necessidade, consistência e verificabilidade.
\end{itemize}

\subsection*{Parte 3: Processo de requisitos (25 min)}
\begin{itemize}[leftmargin=*,itemsep=0.35em]
  \item Elicitação: como coletar necessidades (entrevista, observação, documento).
  \item Análise: conflitos, prioridades, regras e dependências.
  \item Especificação: backlog, histórias de usuário e casos de uso.
  \item Validação: revisão com stakeholders e critérios de aceitação.
\end{itemize}

\subsection*{Parte 4: Atividade prática (35 min)}
\begin{itemize}[leftmargin=*,itemsep=0.35em]
  \item Dividir em grupos (3--4).
  \item Cada grupo:
  \begin{itemize}[leftmargin=*,itemsep=0.25em]
    \item Escreve 3 requisitos funcionais (formato história de usuário).
    \item Escreve 2 requisitos não funcionais (ex.: desempenho, segurança).
    \item Para cada requisito funcional, define 2 critérios de aceitação testáveis.
  \end{itemize}
  \item Revisão rápida: um requisito por grupo no quadro; melhorar clareza e verificabilidade.
\end{itemize}

\subsection*{Parte 5: Fechamento e ponte para prototipagem (20 min)}
\begin{itemize}[leftmargin=*,itemsep=0.35em]
  \item Mostrar como requisitos guiam protótipos e testes.
  \item Reforçar: critérios de aceitação são ``ponte'' para casos de teste.
  \item Anunciar Aula 5: validar interface e fluxo cedo para reduzir retrabalho.
\end{itemize}

\section{Avaliação}
\begin{itemize}[leftmargin=*,itemsep=0.25em]
  \item Coerência e completude do conjunto de requisitos produzido.
  \item Qualidade dos critérios de aceitação (observáveis e testáveis).
  \item Participação e melhoria após feedback.
\end{itemize}

\section{Referências}
\begin{itemize}[leftmargin=*,itemsep=0.25em]
  \item \href{../99-Materiais/Engenharia%20de%20software%20projetos%20e%20processos%20by%20Wilson%20de%20P%C3%A1dua%20Paula%20Filho%20.epub.pdf}{PAULA FILHO, Wilson de Pádua. Engenharia de software: projetos e processos. 4. ed. 2019.}
\end{itemize}

\end{document}
